En base a los experimentos realizados, se concluye que tanto en algoritmos de ordenamiento como en multiplicación de matrices, el desempeño relativo de los algoritmos está fuertemente influenciado por la naturaleza de los datos de entrada. En ordenamiento, aunque todos los algoritmos escalan de manera similar, existe una clara jerarquía de eficiencia con \texttt{std::sort} como el algoritmo óptimo, seguido de Quick Sort y Merge Sort; sin embargo, el comportamiento difiere según el orden inicial de los datos, pudiendo Merge Sort ser significativamente más lento en datos ya ordenados ascendentemente. En multiplicación de matrices, el algoritmo de Strassen confirma su ventaja asintótica para matrices grandes, reduciendo el tiempo de cómputo en comparación con el algoritmo clásico, validando así la importancia de la selección de algoritmos eficientes en operaciones computacionalmente intensivas.
